% ============================================================================
% DROP-IN DRAFT: action-space dissection of the value-based nulls.
% Locked regardless of the MPE outcome. Place in the Generalization section
% (after the cross-algorithm subsection) or as its own short subsection.
% Replace bracketed [...] with final analyzer numbers / CIs before submission.
% ============================================================================

\subsection{Paradigm or action space? A controlled dissection}
\label{sec:action_space_dissection}

Extending the cross-algorithm test to value-based learners raised a sharper
question than replication alone. Value-based methods require a discrete action
space, so we discretized each node's continuous base-stock level into $K{=}11$
bins over $[0,500]$ and trained two value-based learners under the identical
reward pipeline, $w$-grid, $\lambda$-sweep, and collapse label used throughout:
an independent value learner (IQL) and a centralized monotonic-mixing learner
(QMIX). Both produced a \emph{flat, $w$-independent} collapse rate
(QMIX ${\approx}18\%$, IQL ${\approx}0\text{--}15\%$) with no fittable
$w_{\mathrm{crit}}$ in any condition, null included. Mean base-stock parked at a
stable interior level irrespective of $w$ (IQL: $\overline{S}\approx[155,157,152,141,151]$
across the grid), and refining the grid to $K{=}41$ or extending the training
budget did not recover a boundary ($K{=}41$ gave $\overline{S}$ identical at
$w{=}0.5$ and $w{=}0.95$).

Two factors separate these learners from the continuous PPO and MAPPO runs that
yield our headline boundary: the \emph{learning paradigm} (on-policy policy
gradient vs.\ off-policy value-based) and the \emph{action representation}
(continuous vs.\ discrete). We isolate them with a control that varies only the
latter: PPO---the learner that produces a clean loss-aversion boundary under
continuous control---trained on the identical $K{=}11$ discrete action shell,
with every other component held fixed. It too goes flat
(Table~\ref{tab:action_dissection}): mean base-stock overlaps across $w$
($\overline{S}\in[100,200]$ at both $w{=}0.5$ and $w{=}0.95$), profit stays
healthy and $w$-independent ($\approx[438,513]$), and no collapse regime
appears.

This control is decisive for the attribution. Holding the learner fixed at PPO,
\emph{discretization alone is sufficient to erase the boundary}; the QMIX and IQL
nulls therefore require no appeal to the value-based paradigm. We make only this
claim, and not its converse: because no continuous-action value-based learner was
run, we do not assert that a value-based method would \emph{preserve} the
boundary given continuous actions---only that discreteness is sufficient to
account for the observed nulls.

\begin{table}[t]
\centering
\small
\begin{tabular}{lll}
\toprule
Learner & Action space & Boundary? \\
\midrule
PPO / MAPPO              & continuous   & yes (Sec.~\ref{sec:results}) \\
QMIX (centralized value) & discrete $K{=}11$ & no \\
IQL (independent value)  & discrete $K{=}11$ & no \\
PPO (control)            & discrete $K{=}11$ & \textbf{no} \\
\bottomrule
\end{tabular}
\caption{Action-space dissection. The control (last row) changes only the action
representation while holding the learner fixed at PPO, isolating discretization
---not the value-based paradigm---as the cause of the value-based nulls.}
\label{tab:action_dissection}
\end{table}

The mechanism of Sec.~\ref{sec:mechanism} reads this result as a prediction
confirmed rather than a failure. The boundary is located there at a
signal-to-noise crossing of the \emph{continuous} policy gradient: the $w$-blend
moves the \emph{value} of the optimal base-stock level, not \emph{which} of the
coarse bins is the arg-max. A discrete arg-max policy therefore snaps to a stable
interior bin and never visits the near-threshold $(s,S)$ region where collapse
occurs, so the mechanism predicts that coarse discretization should remove the
boundary---which the control confirms. We accordingly scope the phenomenon as one
of \emph{continuous-control} cooperative MARL, replicating across environments and
across continuous policy-gradient learners (PPO, MAPPO), and we report the
discretized regime as a controlled boundary condition rather than a competing
result.
